\begin{opinion}
    En este trabajo se aborda un problema clásico de la Ciencia de la computación, el problema de la satisfacibilidad (SAT), desde una perspectiva que vincula la teoría de lenguajes formales con la complejidad computacional. Se propone una construcción que utiliza autómatas finitos para representar las interpretaciones que satisfacen cada cláusula de manera independiente, para luego intersectarlas con un lenguaje que fuerza la consistencia de la interpretación global. Este enfoque permite reformular el problema de satisfacibilidad como el problema del vacío para un lenguaje específico.

    Un resultado particularmente relevante de esta tesis es el estudio crítico de las gramáticas libres de contexto múltiples paralelas (pMCFG). Jossué no solo aplica formalismos existentes, sino que identifica un error en el procedimiento de intersección reportado en la literatura, lo que constituye un hallazgo científico significativo. Además, demuestra que el problema de intersección entre una pMCFG y un lenguaje regular es NP-duro, contribuyendo así a una mejor comprensión de los límites de este formalismo.

    Como podrán comprobar las personas que lean este documento, Jossué tuvo que estudiar contenidos que no forman parte de su plan de estudios, incursionando en temas avanzados de teoría de lenguajes, autómatas y complejidad, y lo hizo de manera creativa y rigurosa. También demostró capacidad para identificar problemas abiertos o errores en la literatura existente y proponer soluciones originales.

    Estoy muy satisfecho con la independencia demostrada por Jossué durante todo el proceso de investigación, la calidad de los resultados obtenidos, así como la profundidad del análisis teórico presentado.

    Creo que estamos en presencia de un trabajo excelente, realizado por un excelente científico de la computación.

    \begin{flushright}
    \underline{\hspace{6.5cm}}\\
    MSc. Fernando Raul Rodriguez Flores\\

    Facultad de Matemática y Computación
    
    Universidad de la Habana

    Junio, 2026
    \end{flushright}
\end{opinion}